% ------------------------------------------------------------
% Appendix C — Derived-Geometry Glossary for Engineers
% ------------------------------------------------------------

\chapter{Derived-Geometry Glossary for Engineers}

The purpose of this appendix is to provide a coherent and mathematically rigorous exposition of the fundamental objects of derived geometry as they relate to delta–sigma modulation and the broader RSVP field-theoretic framework. Although these notions arise from algebraic geometry, higher category theory, and homotopical methods, they have concrete meaning when interpreted through the geometry of quantization, entropy gradients, and feedback control. The key principle is that quantization produces singular behaviors and discontinuous projections that cannot be represented fully using classical smooth manifolds. Derived geometry provides an appropriate language for capturing these features by encoding both the visible state manifold and the hidden infinitesimal structures that emerge at quantizer boundaries.

In what follows, each concept is defined with formal precision, followed by a translation into the intuition relevant for delta–sigma modulators. Because this book reinterprets measurement, quantization, and feedback control as interactions between scalar, vector, and entropy fields, the derived-geometric viewpoint becomes the natural generalization of classical linear-systems theory. This appendix therefore functions not only as a glossary but also as a conceptual bridge between engineering practice and modern geometric frameworks.

\section{C.1 Derived Manifolds and Infinitesimal Structure}

A classical manifold is defined purely by its local coordinate structure. Derived manifolds extend this notion by encoding not only the visible points but also the infinitesimal directions in which classical smoothness fails. Formally a derived manifold is a space equipped with a sheaf of simplicial commutative rings, such that the truncation of this sheaf recovers the smooth function sheaf of a classical manifold. The higher homotopy groups of this sheaf encode the non-smooth or obstructed directions.

In the context of delta–sigma modulation, derived manifolds arise naturally because quantizer boundaries impose discontinuous projection maps. At these boundaries the tangent space is not well-defined in the classical sense. The derived tangent complex records the directions along which the system would have evolved if not for the discontinuous collapse imposed by the quantizer. These invisible directions correspond to the residual error dynamics that shape the noise spectrum. The derived manifold therefore captures the latent geometry of the modulator that classical manifolds erase.

\section{C.2 Derived Stacks and Families of Modulators}

Whereas derived manifolds describe fixed configuration spaces, derived stacks describe entire families of objects varying over a base. A derived stack is a functor from derived rings to simplicial sets that satisfies descent conditions. In practice, a derived stack allows one to parameterize spaces of solutions, control architectures, or feedback configurations in a way that accommodates singularities and higher symmetries.

A delta–sigma modulator with tunable parameters such as loop filters, quantizer levels, or oversampling ratio forms such a family. The derived stack of modulators encodes not only the visible mapping from parameters to states but also the infinitesimal degeneracies and local obstructions that may appear when filter poles approach instability, quantizer boundaries collide, or oversampling ratios are changed. Classical parameter spaces such as \(\mathbb{R}^{n}\) cannot describe these phenomena adequately; derived stacks record the higher-order behavior needed to fully represent how modulators deform under parameter changes.

\section{C.3 Homotopy Limits and Feedback Composition}

In derived geometry the appropriate notion of a limit is the homotopy limit. Unlike ordinary limits, homotopy limits remember the higher morphisms that arise when diagrams fail to commute strictly. The homotopy limit of a diagram represents the universal object that coherently satisfies all morphisms up to homotopy.

Feedback systems in engineering rarely commute strictly; delays, nonlinearity, quantization, and approximate control break exact commutativity. When a continuous-time system is sampled and passed through a quantizer, the loop no longer corresponds to a strict pullback diagram but instead closes only up to homotopy. Thus the modulator loop is properly represented as a homotopy limit. The homotopy data encode the spectral distribution of the quantization error, the mismatch between predicted and actual states, and the residual curvature introduced by the discrete collapse.

This identification provides a conceptual explanation for why the block diagram equations of delta–sigma modulators must be interpreted carefully when discussing stability. Classical limits ignore the mismatch introduced by quantizer impulses. The homotopy limit integrates this mismatch into the very definition of the closed loop.

\section{C.4 Fiber Products and Symbolic Collapse}

The fiber product is the classical construction used to impose constraints on manifolds. In derived geometry the fiber product is homotopy-corrected, producing a derived fiber product that preserves invisible infinitesimal structure. At a quantizer boundary the modulator state is projected onto a discrete set of symbols. This projection corresponds mathematically to intersecting the continuous state manifold with the discrete symbol manifold. The intersection is never transverse; the continuous manifold generically intersects the discrete set in a singular manner.

Thus the quantizer action is precisely a derived fiber product. The homotopical correction term records the lost tangent directions. These directions correspond to the error dynamics that propagate through the loop filter. The NTF zeros arise from the annihilation of certain derived tangent directions. Idle tones arise when the derived tangent complex aligns with periodic symbolic strata, producing persistent cycles. The derived fiber product therefore provides a rigorous geometric account of quantization behavior.

\section{C.5 Shifted Symplectic Structures and Energy–Error Duality}

The theory of shifted symplectic structures generalizes classical symplectic geometry to higher stacks. A \(k\)-shifted symplectic structure is a closed nondegenerate two-form of degree \(k\) on a derived stack. These structures govern the behavior of field theories, topological invariants, and moduli spaces.

In delta–sigma modulation the error dynamics can be interpreted as a discrete field theory. The error acts as a canonical coordinate, while the feedback acts as a conjugate momentum. The modulator thus possesses a naturally shifted symplectic structure that reflects the duality between energy storage in the loop filter and error correction through feedback. The shift \(k = -1\) corresponds to the fact that the quantizer projects onto zero-dimensional fibers, reducing phase-space dimension by one.

The shifted symplectic viewpoint reveals that delta–sigma modulators are not merely dissipative systems. They possess a regulated exchange between energy-like quantities stored in integrators and entropy-like quantities discharged through quantization. This duality mirrors the RSVP decomposition into scalar potentials, vector flows, and entropy fields.

\section{C.6 Derived Cotangent Complex and Higher-Order Delta–Sigma Perturbations}

The cotangent complex is the derived replacement for the cotangent bundle. It encodes all infinitesimal deformations of a derived space and includes higher homotopy information. In modulators the cotangent complex describes how infinitesimal perturbations evolve under both the vector field and the quantizer collapse. Unlike the classical Jacobian, the derived cotangent complex incorporates the hidden modes that re-emerge after quantization through the loop filter.

The structure of the cotangent complex determines the sensitivity of the modulator to dither, input slope, and loop-filter misalignment. It also predicts the onset of chaotic or quasi-periodic regimes in higher-order modulators. When the cotangent complex acquires nontrivial higher homology, the system exhibits persistent oscillations because the derived tangent directions are not fully damped by entropy descent. Thus the cotangent complex provides a precise mathematical framework for understanding idle tones, bifurcations, and harmonic distortion.

\section{C.7 Derived Mapping Spaces and Time-Varying Quantizers}

A mapping space in derived geometry is itself a derived stack that parameterizes all maps between two stacks. Time-varying systems, such as delta–sigma modulators with adaptive quantizers or variable loop filters, correspond to paths in a derived mapping space. The geometry of this mapping space determines the smoothness of parameter evolution and the stability of adaptive control laws.

When the quantizer thresholds vary slowly, the system traverses a smooth path in the mapping space, and stability can be maintained. Rapid variation corresponds to moving through nontrivial homotopy classes, which leads to symbolic discontinuities and transient bursts of error. The mapping-space perspective therefore yields a formal method for analyzing the stability of adaptive delta–sigma modulators, especially those used in machine-learning-driven feedback architectures.

\section{C.8 Integration of Derived Geometry with the RSVP Framework}

The field-theoretic interpretation developed throughout this book naturally incorporates derived geometry. RSVP theory treats systems as interactions between scalar potentials, vector flows, and entropy fields. Derived geometry augments this picture by encoding the discontinuities introduced by quantization and the infinitesimal memory stored in the tangent complex of the derived state manifold.

The entropy field corresponds to the negative curvature directions in the derived tangent complex. The vector flow corresponds to the derived pushforward of tangent vectors. Quantization corresponds to derived symbolic collapse. Stability corresponds to the existence of a shifted symplectic potential that bounds the derived energy–entropy exchange.

Thus derived geometry provides the mathematical infrastructure necessary to integrate delta–sigma theory into RSVP dynamics. It reveals that quantization is not merely a nonlinear artifact but a fundamental geometric operation that shapes the informational topology of measurement.
